\subsection{The distinction between absolute and relative Theorem A}

\begin{claim}{B.1}{Proved}{Absolute Theorem A for every derivator}
For every derivator $\D$, the class $\Wlc$ is weakly saturated,
contains the categories with a terminal object, and satisfies the
absolute Theorem A: if $u:A\to B$ and every $A/b\to e$ belongs to
$\Wlc$, then $u\in\Wlc$. The constant-coefficient class $\Wc$
is a fundamental localizer, including the relative Theorem A.
The class $\Wone$ need not be a localizer, by A.5.
\end{claim}
\begin{proof}
Weak saturation is A.4. A category with a terminal object has
contractible nerve, so A.2 applies. To check the first Theorem A,
we use covariant coefficients $\mathbb E(I)=\D(I^{\op})$.
Passage to opposites in A.7 identifies its classes with those of
$\D$. Each transition $A/b\to A/b'$ is a $\Wlc$-equivalence:
its two projections to $e$ are, and two-out-of-three applies.
For $F\in\LC_{\mathbb E}(A)$, pointwise left Kan extension gives
\[
 (u_!F)_b=\colim_{A/b}F|_{A/b}.
 \tag{B.1.1}
\]
Here and below the displayed colimit denotes the derivator Kan
extension, an object of $\mathbb E(e)$. A transition $b\to b'$
induces an isomorphism between these colimits: A.2 applied to the
transition category and the restricted locally constant coefficient
is exactly this comparison. Hence $u_!F$ is locally constant.

For each object $(a,\alpha)$ of $A/b$, restriction from
$\LC_{\mathbb E}(A/b)$ to $\mathbb E(e)$ at that object is an
inverse to the constant-diagram equivalence. Its value-to-colimit
map is therefore an isomorphism. In particular, at
$(a,\id_{u(a)})$, the unit $F\to u^*u_!F$ is an isomorphism.
For $G\in\LC_{\mathbb E}(B)$ the restriction of $u^*G$ to
$A/b$ is canonically isomorphic to the constant value $G_b$,
using the arrows $u(a)\to b$. The counit $u_!u^*G\to G$ is
therefore an isomorphism at $b$. Point evaluations detect both
isomorphisms. These are the specified unit and counit, so they
prove that $u^*$ is an equivalence on locally constant objects.
This is the argument of \citet[Proposition~1.27]{Cisinski2009LocallyConstant},
whose Kan-extension argument applies here using A.6.

For the last assertion we can use
\citet[Proposition~3.7]{Cisinski2008}; its localizer axiom is
explicitly the relative one. To specify why it concerns our class,
Yoneda identifies (A.1.1) with invertibility of
$p_{B*}p_B^*\to p_{A*}p_A^*$. For a triangle
$A\xrightarrow{u}B\xrightarrow{\beta}S$, with
$\alpha=\beta u$, the pointwise formula identifies the value at
$s$ of $\beta_*p_B^*\to\alpha_*p_A^*$ with the constant
cohomology comparison for $A/s\to B/s$. If all these comparisons
are isomorphisms, Der~2 and then $p_{S*}$ make the comparison on
global constant cohomology an isomorphism. This proves relative
Theorem A with all maps typed in $\D(S)$ and $\D(e)$.
\end{proof}

\begin{claim}{B.2}{Proved}{Relative Theorem A with an infinity-category enhancement}
Let $\mathcal C$ be a complete and cocomplete infinity-category,
with the universe bounds of the conventions. Then
$\W^{\mathrm{lc}}_{\mathcal C}$ is a fundamental localizer.
In particular, for a morphism $u:A\to B$ over a small category
$S$, the hypotheses
\[
 (A/s\longrightarrow B/s)\in\W^{\mathrm{lc}}_{\mathcal C}
 \quad(s\in S)
 \tag{B.2.1}
\]
imply that the specified restriction
$\Fun(BB,\mathcal C)\to\Fun(BA,\mathcal C)$ is an
equivalence of infinity-categories.
\end{claim}
\begin{proof}
The hypotheses initially concern homotopy categories, but A.8
upgrades each specified restriction to an infinity-categorical
equivalence, using its Kan adjunctions. Consider the covariant
diagram $s\mapsto A/s$. Its Grothendieck construction has objects
$(s,a,p(a)\to s)$, where $p:A\to S$. The projection
$P:\int_S(A/s)\to A$ has the section
$J(a)=(p(a),a,\id)$, and $J\dashv P$. Thus $BP$ is a weak
homotopy equivalence. The Grothendieck-construction formula for
homotopy colimits, in the form
\citet[Theorem~1.13 and Remark~1.14]{Cisinski2008}, gives natural
equivalences
\[
 BA\simeq\hocolim_{s\in S}B(A/s),\qquad
 BB\simeq\hocolim_{s\in S}B(B/s).
 \tag{B.2.2}
\]
The space underlying this homotopy colimit is also the colimit of
the same diagram in infinity-categories: the inclusion of spaces
is left adjoint to the maximal-groupoid functor. Cartesian closure
then gives, in the larger universe,
\[
 \Fun(BA,\mathcal C)\simeq
 \lim_{s\in S^{\op}}\Fun(B(A/s),\mathcal C),
 \tag{B.2.3}
\]
and the corresponding formula for $B$. The diagram of restriction
functors is coherently the one induced by $u$; (B.2.1) makes it an
objectwise equivalence. Its limit is an equivalence, proving the
claim about the specified map. A.4 and A.2 supply the remaining
localizer axioms. If $S$ is empty, both $A$ and $B$ are empty,
so the same formula and conclusion hold.
\end{proof}

\begin{claim}{B.3}{Proved}{Relative Theorem A for an arbitrary derivator}
Let $\D$ satisfy Der~1--4 in the presheaf convention of A.1, on all
$\U$-small categories. Let $A\xrightarrow{u}B\xrightarrow{q}S$ be
functors of $\U$-small categories and put $p=qu$. If, for every
$s\in S$, the specified restriction
\[
 (u/s)^*:\LC_\D(B/s)\longrightarrow\LC_\D(A/s)
 \tag{B.3.1}
\]
is an equivalence, then $u^*:\LC_\D(B)\to\LC_\D(A)$ is an
equivalence. If every restriction in (B.3.1) is fully faithful,
then $u^*$ is fully faithful. Both conclusions also hold for a map
of Grothendieck opfibrations preserving cocartesian arrows when the
corresponding restriction on every strict fiber is, respectively,
an equivalence or fully faithful. Consequently $\Wlc$ is a
fundamental localizer for every such derivator, while $\Wone$
satisfies relative Theorem A despite its failure of
two-out-of-three in A.5. The statement concerns the ordinary categories
$\LC_\D(-)$ themselves; neither an infinity-categorical presentation
nor strongness of $\D$ is required.
\end{claim}
\begin{proof}
We first establish the partial localizations needed to make descent
coherent. If $T$ is small and $R\subseteq T$ is a subcategory, let
$\D(T;R)$ be the full subcategory of $\D(T)$ on objects whose
underlying presheaf inverts every arrow of $R$. Both adjoints to
\[
 j_{T,R}:\D(T;R)\hookrightarrow\D(T)
 \tag{B.3.2}
\]
exist, and they are the components of adjunctions of shifted
derivators. Here and below ``shifted'' means that, at an auxiliary
shape $J$, the same condition is imposed on the $R$ arrows in
$T\times J$.

For completeness, this is the arbitrary-subcategory form of the
argument in \citet[Theorem~3.6, Lemma~3.9, Proposition~3.10, and
Theorem~3.13]{Cisinski2009LocallyConstant}. Write $C=T^{\op}$ and
$R'=R^{\op}$. The universal derivator of simplicial presheaves on
$C$ has a localization
\[
 \gamma:\Ho([C^{\op},\mathrm{sSet}])
 \rightleftarrows
 \Ho(L_{R'}[C^{\op},\mathrm{sSet}]):i,
 \qquad\gamma\dashv i.
 \tag{B.3.3}
\]
The right adjoint $i$ is fully faithful. It preserves homotopy
colimits: a pointwise homotopy colimit of diagrams which invert
$R'$ again inverts $R'$, since a homotopy colimit sends an
objectwise weak equivalence of diagrams of spaces to a weak
equivalence. This argument works at every auxiliary shape. Thus both
$\gamma$ and $i$ are cocontinuous morphisms of derivators.
Cisinski's Theorem~3.6, which is explicitly stated for an arbitrary
subcategory $R'$ and every derivator, identifies precomposition by
$\gamma$ on cocontinuous morphisms into $\D$ with (B.3.2).
Precomposition by $i$ supplies its left adjoint. The same construction
with the opposite derivator and opposite subcategory supplies its
right adjoint. The internal-Hom argument of Theorem~3.13 gives these
as adjunctions of derivators, because the argument just given applies
to every shift. All the categories and sets of arrows used in the
localization are $\U$-small; the simplicial-presheaf categories are
taken in the ambient larger universe. This uses the universal
property of the auxiliary model derivator, not a presentation of
$\D$ by a model category.
There is no hidden strongness hypothesis in this universal property:
the cited foundational paper defines a two-sided derivator using
Der~1--4, and Proposition~3.12 makes it eligible for Corollary~3.26
on every small shape
\citep[Definitions~0.1--0.2, Proposition~3.12, and
Corollary~3.26]{Cisinski2008}.

We use the right adjoint of (B.3.2), denoted $Q_{T,R}$. Suppose that
$\pi:T\to S$ is a Grothendieck opfibration and let $R$ consist of
its vertical arrows, namely those sent to identities. Write
$\D(T)^{\mathrm{vert}}=\D(T;R)$ and $Q_T=Q_{T,R}$. Its objects
are exactly the diagrams locally constant on every fiber. For the
fiber inclusion $i_s:T_s\hookrightarrow T$, the canonical mate is
an isomorphism
\[
 i_s^*Q_T\ \cong\ Q_{T_s}i_s^*.
 \tag{B.3.4}
\]
We verify the required base change rather than assume descent.
The left Kan extension $i_{s!}$ sends locally constant objects on
$T_s$ to vertically locally constant objects on $T$. Indeed, at
$x\in T_t$, the comma category $x\backslash T_s$ is the disjoint
union, over $\alpha:t\to s$, of categories with initial object
the chosen cocartesian lift $x\to\alpha_!x$. Der~4, in presheaf
variance, therefore gives
\[
 x^*i_{s!}F\ \cong\
 \coprod_{\alpha\in\Hom_S(t,s)}(\alpha_!x)^*F.
 \tag{B.3.5}
\]
For a vertical arrow $x\to x'$, the comparison in (B.3.5) is the
coproduct of the comparisons induced by the vertical arrows
$\alpha_!x\to\alpha_!x'$. They are isomorphisms when $F$ is
locally constant. The assertion follows, including empty comma
categories and infinite coproducts. The equality of the two
composites ``include, then $i_{s!}$'' and ``$i_{s!}$ on the local
subcategories, then include'' now gives (B.3.4) by taking right
adjoints. Consequently (B.3.4) is compatible with the units and
counits, and the argument also proves it at every shift.

Consider a map $v:T\to T'$ of Grothendieck opfibrations over $S$
preserving cocartesian arrows. For $x'\in T'_s$, the inclusion
\[
 T_s/x'\hookrightarrow T/x'
 \tag{B.3.6}
\]
has a left adjoint. It sends $(x,\beta:v(x)\to x')$, with
$\pi(x)=t$ and $\pi'(\beta)=\alpha:t\to s$, to the object
$(\alpha_!x,\overline\beta)$ obtained by factoring $\beta$
through the cocartesian transport of $v(x)$. The universal property
of a cocartesian arrow gives the adjunction. If $r\dashv i$ is
this adjunction, presheaf variance gives $r_*\cong i^*$; composing
right Kan extensions to the point shows that restriction along
(B.3.6) computes the same homotopy limit. Der~4 therefore gives the
canonical isomorphism
\[
 i_s'^*v_*\ \cong\ v_{s*}i_s^*.
 \tag{B.3.7}
\]
This is the usual base change for a map of opfibrations, proved here
with the variance required for $\D$.

Suppose now that every $v_s^*:\LC_\D(T'_s)\to\LC_\D(T_s)$
is an equivalence. On vertically locally constant diagrams we have
an adjunction
\[
 v^*:\D(T')^{\mathrm{vert}}
 \rightleftarrows\D(T)^{\mathrm{vert}}:Q_{T'}v_*.
 \tag{B.3.8}
\]
By (B.3.4) and (B.3.7), its unit and counit restrict on every fiber
to those of
$v_s^*\dashv Q_{T'_s}v_{s*}$. Those units and counits are
isomorphisms by hypothesis. Every point lies in a fiber, so Der~2
makes the unit and counit of (B.3.8) isomorphisms. Thus (B.3.8) is
an equivalence.

If the fiber restrictions are only fully faithful, the same argument
applied to the unit alone proves that $v^*$ in (B.3.8) is fully
faithful. Its restriction to the full subcategories of fully locally
constant diagrams is therefore fully faithful as well. This proves
the asserted fiberwise full-faithfulness conclusion without an
essential-surjectivity or reflection argument.

It remains to restrict this equivalence to diagrams locally constant
on all arrows. Restriction $v^*$ preserves this condition. It also
reflects it among vertically locally constant diagrams. To see this,
let $N\in\D(T')^{\mathrm{vert}}$ and suppose $v^*N$ is locally
constant. For $\alpha:s\to t$, the cocartesian lifts define a
natural transformation between the corresponding fiber inclusions,
and hence a morphism
\[
 \alpha_!^*(i_t'^*N)\longrightarrow i_s'^*N
 \quad\hbox{in }\LC_\D(T'_s).
 \tag{B.3.9}
\]
Its image under $v_s^*$ is an isomorphism, since it is the transport
comparison for $v^*N$. The equivalence $v_s^*$ is conservative, so
(B.3.9) is an isomorphism. Every arrow of $T'$ factors as a
cocartesian arrow followed by a vertical arrow. Equation (B.3.9)
and vertical local constancy therefore show that $N$ is locally
constant. Essential surjectivity in (B.3.8), followed by this
reflection, gives essential surjectivity on fully locally constant
diagrams; full faithfulness is inherited from (B.3.8). We have proved
that $v^*:\LC_\D(T')\to\LC_\D(T)$ is an equivalence.

Finally replace $p:A\to S$ by the Grothendieck opfibration
\[
 \widetilde A=\int_{s\in S}(A/s)\longrightarrow S.
 \tag{B.3.10}
\]
An object is $(a,s,\alpha:p(a)\to s)$, and a morphism is a pair
$(h:a\to a',k:s\to s')$ with $k\alpha=\alpha'p(h)$.
Transport along $k$ replaces $\alpha$ by $k\alpha$. The forgetful
functor $r_A:\widetilde A\to A$ has the left adjoint
$a\mapsto(a,p(a),\id)$, with $r_Aj_A=\id$ and
$j_Ar_A\Rightarrow\id$. The latter natural transformation induces
an isomorphism on every locally constant object, by Der~2.
Consequently $r_A^*$ and $j_A^*$ are mutually inverse equivalences
on locally constant categories. The same holds for $B$. The functor
$u$ induces a transport-preserving map
$\widetilde u:\widetilde A\to\widetilde B$ whose fiber at $s$
is precisely $u/s$. Apply (B.3.8)--(B.3.9) to $\widetilde u$ and
then use the forgetful equivalences. For full faithfulness use the
unit-only version of (B.3.8) just proved. This proves both relative
Theorem A assertions.

The unit, counit, and base-change proof is the ordinary-derivator
specialization of the proof of
\citet[Main Theorem~2.5.5(2), pp.~52--54]{EF_Hoermann2015}.
The stable compact-generation assumptions in that fibered theorem
provide its partial Cartesian projectors. In the present
nonfibered situation, (B.3.2)--(B.3.3) provide those projectors for
every derivator, so those assumptions are unnecessary. Weak
saturation was proved in A.4. A category with a terminal object is
$\Wlc$-aspherical by the same adjoint-section argument as for
$r_A$, using its projection to the point and the terminal section.
Together with the relative theorem just proved, these are precisely
the three fundamental-localizer axioms.
\end{proof}


\begin{claim}{B.4}{Proved}{Accessibility for all computed coefficient theories}
For every $\U$-presentable enhanced coefficient category, and in
particular for the derivators of sets, groupoids, $r$-types ($r\geq0$), spaces,
$D(k)$ for a small associative unital ring $k$, spectra, and $L_E\Sp$
for a small spectrum $E$, the constant and $2$-complete classes are
set-generated fundamental localizers. They also satisfy the literal
accessibility condition Loc~(9) of Chapter XVI ms.~75--79: as full
subcategories of $\Cat^{[1]}$, they are generated under sufficiently
filtered colimits by a small family of bounded-cardinality arrows.
For the listed concrete coefficient theories the intermediate
classes are accessible full subcategories of
$\Cat^{[1]}$ as well; this does not turn them into fundamental
localizers when two-out-of-three fails. No large-cardinal principle
is used for these coefficient theories.
\end{claim}
\begin{proof}
We give the bounded-subfunctor construction, so that the two meanings
of accessibility are both verified. First fix $k\ne0$ and choose an
uncountable regular cardinal $\kappa\in\U$ larger than $|k|$ and
$\aleph_0$. Size of a category means its number of arrows, including
identities. Put $Q=\operatorname{Ex}^{\infty}$. The functor $Q$
preserves monomorphisms and filtered colimits. Every simplex of $QK$
occurs at a finite $\operatorname{Ex}$-stage and is a map from a
finite subdivided simplex; consequently it depends on finitely many
simplices of $K$. If a category has fewer than $\kappa$ arrows, its
$Q$-nerve has fewer than $\kappa$ simplices in total. For a vertex
$a$ of a Kan complex $K$, use the based-loop model
\[
 \Omega_aK=\{(a,a)\}\times_{K\times K}K^{\Delta^1}.
\]
The endpoint map is a Kan fibration. Each loop simplex uses only the
finite simplicial set $\Delta^n\times\Delta^1$, so this construction
also has finite support and commutes with filtered colimits of
pointed Kan complexes.

Let $f:A\to B$ belong to $W_{D(k)}^{\mathrm{lc}}$, and start with
any subfunctor $g_0:A_0\to B_0$ having fewer than $\kappa$ arrows.
Here $A_0,B_0$ are subcategories and $f(A_0)\subseteq B_0$.
We enlarge it inside $f$ using the criterion of C.7. For every
vertex $b\in B_0$, choose a vertex $a\in A$ and a finite zigzag
from $f(a)$ to $b$. For every pair $a,a'\in A_0$ whose images are
connected in $B_0$, choose a finite zigzag joining them in $A$.
Ambient bijectivity on components supplies these choices.

For every $a\in A_0$, every $n\geq0$, and every cycle
$z\in C_n(\Omega_{f(a)}QN B_0;k)$, choose an ambient cycle
$w\in C_n(\Omega_aQN A;k)$ and an ambient chain $b$ with
$f_*(w)-z=db$. Ambient surjectivity on loop homology supplies them.
For every cycle $w\in C_n(\Omega_aQN A_0;k)$ and every chain
$b\in C_{n+1}(\Omega_{f(a)}QN B_0;k)$ with $db=f_*(w)$, choose
an ambient chain $c$ with $dc=w$; ambient injectivity supplies it.
Each selected chain and zigzag has finite support. Pull the
supporting $Q$-nerve data back to finitely many original arrows of
$A$ and $B$, add those arrows and their images under $f$, and
generate subcategories by identities and finite composition.
There are fewer than $\kappa$ tests and witnesses: all relevant
chain groups have fewer than $\kappa$ elements because
$|k|<\kappa$. Thus the enlarged subfunctor $g_1$ still has size
less than $\kappa$.

Repeat this construction countably many times. Its union $g_\infty$
has size less than $\kappa$ by regular uncountability. Every
component relation in the union has a finite zigzag witness and
occurs at a finite stage. Every cycle and every bounding chain
likewise occurs at a finite stage, because its support is finite
and $QN$ and the loop model commute with filtered unions. At the
next stage the required lift or bounding chain was added. Thus
$g_\infty$ is a component bijection and its based-loop maps are
$k$-homology equivalences. C.7 gives
$g_\infty\in W_{D(k)}^{\mathrm{lc}}$.

Let $P_f$ be the poset of all size-$<\kappa$ subfunctors of $f$
lying in this class. Given fewer than $\kappa$ members, take their
union, generate subcategories, and apply the construction just
made. This gives an upper bound in $P_f$, so $P_f$ is
$\kappa$-filtered. Every original arrow lies in some member,
by starting the construction with that arrow. Therefore the
specified inclusions give
\[
 f=\colim_{g\in P_f}g\quad\text{in }\Cat^{[1]}.
\]
This also checks the enlargement form of the manuscript's Loc~(9),
with inclusions of source into source and target into target.
The same finite-support tests show closure of the class under
filtered colimits: $QN$, based loops with a chosen vertex, and
homology commute with them; bijectivity on components is detected
by finite zigzags. A category with fewer than $\kappa$ arrows is
$\kappa$-presentable in $\Cat$, and similarly for a functor in
$\Cat^{[1]}$: all its generators and composition equations form
a set of size $<\kappa$. These small weak arrows remain
$\kappa$-presentable in the full subcategory because its filtered
colimits are computed in the ambient arrow category. Choosing
representatives of their isomorphism classes therefore proves
accessibility as a full subcategory.

For set-generation as a fundamental localizer, let $S$ be that
small set of representatives and let $\langle S\rangle$ be its
generated fundamental localizer. B.2 gives
$\langle S\rangle\subseteq W_{D(k)}^{\mathrm{lc}}$.
Every $g\in P_f$ belongs to $\langle S\rangle$. Every
fundamental localizer is closed under filtered colimits, by
\citet[Corollary~4.2.22]{Cisinski2006}; hence
$f\in\langle S\rangle$. This proves equality. The cited
corollary alone is not an accessibility theorem; the bounded
construction is the indispensable extra argument.

For $L_E\Sp$, represent $E$ by a small cofibrant sequential
spectrum and choose $\kappa$ also larger than the total number
of simplices in all its levels. Use C.11 and replace loop-chain
cycles and boundaries by representatives and equality witnesses
in the stable homotopy groups of
$E\wedge\Sigma^\infty_+\Omega_aQN A$. A stable class is
represented at a finite spectrum level by a finite simplicial
sphere after levelwise fibrant replacement; an equality or
nullity is witnessed at a further finite level by a finite
homotopy. Thus each witness uses finitely many simplices of
the loop complex. For a size-$<\kappa$ complex there are fewer
than $\kappa$ such representatives and witnesses, over all
integer degrees and spectrum levels. The same countable
enlargement, filtered-union verification, and generated-localizer
argument therefore apply. This requires neither connectivity
of $E$ nor a smashing localization. Spectra are the case
$E=\mathbb S$.

For the constant classes, omit the based-loop construction and
use the ordinary chain or suspension-spectrum homology test
of C.7 or C.11. Their fundamental-localizer assertion is B.1.
For the intermediate classes use the homotopy-fibre test C.2
instead of based loops. A homotopy fibre of a Kan map is
constructed from a path object and a finite pullback, so it
has the same finite-support and filtered-colimit properties.
The resulting full subcategories are accessible, but the
set-generation conclusion as a fundamental localizer is not
asserted for them.

For sets, groupoids, and $r$-types, C.4 reduces the tests to
finitely many homotopy groups and a surjectivity condition
in the next degree. Representatives and homotopies are finite
sphere and disk data in Kan replacements, so the same
enlargement applies with any uncountable regular $\kappa$.
For spaces use every homotopy degree, still a countable family
of finite-support tests. These prove literal accessibility
and, for their constant and $2$-complete fundamental
localizers, set-generation. Zero coefficients give the whole
arrow category and are included by discarding all tests.
\paragraph{All presentable enhanced coefficients.}
We strengthen the concrete argument to every $\U$-presentable
$\infty$-category $\mathcal C$. Choose an uncountable regular
$\kappa\in\U$ for which $\mathcal C$ is
$\kappa$-compactly generated. The category $\PrL_\kappa$ of
$\kappa$-compactly generated presentable categories and
colimit-preserving functors preserving $\kappa$-compact objects
is itself $\U$-presentable:
\citet[Proposition~5.5.7.10]{EF_LurieHTT} identifies it with
$\Cat_\infty^{\mathrm{Rex}(\kappa)}$, and
\citet[Lemma~4.8.4.2]{EF_LurieHA} applies to the
$\U$-small family of representatives of the $\kappa$-small simplicial shapes.
By \citet[Proposition~5.5.7.6 and Notation~5.5.7.7]{EF_LurieHTT},
limits in its dual category $\operatorname{Pr}^{\mathrm R}_\kappa$ are
computed in categories. Consequently
\[
 F_{\mathcal C}:\mathcal S\longrightarrow\PrL_\kappa,
 \qquad X\longmapsto\Fun(X,\mathcal C),\qquad f\longmapsto f_!,
\]
is the copower functor $X\mapsto\colim_X\mathcal C$ and
preserves all small colimits. The natural direction here is
left Kan extension, not restriction; its adjoint is the
specified restriction $f^*$. In particular $F_{\mathcal C}$
is accessible. Restriction is an equivalence exactly when
its adjoint is, so $\W^{\mathrm{lc}}_{\mathcal C}$ is the
inverse image under $\operatorname{Arr}(F_{\mathcal C})$
of the full subcategory of equivalences. For an accessible
category $\mathcal E$, that full subcategory of
$\operatorname{Arr}(\mathcal E)$ is accessible: it is
equivalent to $\mathcal E$ by $e\mapsto\id_e$, and its
inclusion preserves filtered colimits.
\citet[Proposition~5.4.6.6]{EF_LurieHTT}, applied to the
resulting homotopy pullback, proves accessibility of this
arrow class of spaces.

To obtain the literal class of ordinary functors, apply
the same pullback theorem to the accessible functor
\[
 N(\operatorname{Arr}(\Cat_\U))\longrightarrow
       \operatorname{Arr}(\mathcal S),\qquad u\longmapsto Bu.
\]
It preserves filtered colimits because nerves and classifying
spaces do. Since the equivalence condition is full and replete,
this pullback is the nerve of the ordinary full subcategory
of arrows in $\W^{\mathrm{lc}}_{\mathcal C}$. Accessibility
of that nerve is ordinary accessibility: compactness is
detected by its discrete mapping spaces, and a filtered
infinity-category admits a cofinal ordinary filtered poset
\citep[Proposition~5.3.1.18]{EF_LurieHTT}.
Thus a $\U$-small set of functors generates the class under
sufficiently filtered colimits. It is closed under all
filtered colimits by B.2 and Cisinski's Corollary~4.2.22.
Its closure under all filtered colimits is therefore exactly
the class, proving literal Loc~(9). The generated-localizer
argument above proves set-generation as well.

For the constant class, choose a $\U$-small set $G$
generating $\mathcal C$ under colimits. The colimit-preserving
functor $X\mapsto(X\otimes g)_{g\in G}$ from spaces to
$\mathcal C^G$ inverts precisely $\W_{\mathcal C}$ by C.3:
the equivalence test is closed under colimits in the coefficient
object, so it suffices to test the generators. The identical
inverse-image and ordinary-nerve argument applies. This
strengthens all the concrete constant and locally constant
calculations without a large-cardinal hypothesis.
\end{proof}

The order printed on pairs of subcategories in Chapter XVI
ms.~78 reads $X_i\subset Y_i$ and $X_j\subset Y_j$.
It must read $X_i\subset X_j$ and $Y_i\subset Y_j$.
The former expression is not even typed when source and
target categories differ; the proof above uses the latter.

For an arbitrary derivator, accessibility is not proved here.
The bounded tests above depend on a fixed small ring, spectrum,
or truncation theory. Filtered-colimit closure by itself does
not imply set-generation. If one assumes Vop\v enka's
principle, every fundamental localizer is accessible
\citep[Scholie~4.2.17]{Cisinski2006}; that additional
set-theoretic hypothesis is not assumed in this paper.


\begin{claim}{B.5}{Proved}{Sums, components, direct excision, and filtered colimits}
For every derivator, $\Wlc$ satisfies the full disjoint-union
axiom Loc~(4). It satisfies Loc~(5), namely $\Wlc\subseteq\W_0$,
exactly when $\D(e)$ is not equivalent to the terminal category.
For every derivator it also satisfies
the direct Mayer--Vietoris axiom Loc~(6a) and the filtered-colimit
axiom Loc~(8). The class $\Wone$ satisfies Loc~(4), Loc~(6a),
and Loc~(8) for every derivator, despite its possible failure
of two-out-of-three.
For arbitrary $\D$, it satisfies Loc~(5) exactly when $\D(e)$
is not a preorder up to equivalence. For arbitrary $\D$, the constant
class $\Wc$ satisfies Loc~(6a) and Loc~(8); it satisfies Loc~(5)
exactly under the nonthinness condition G.2, and need not satisfy
the reverse implication in Loc~(4).
\end{claim}
\begin{proof}
A.7 proves Loc~(4), and G.1 proves the stated exact component
criterion. For Loc~(6a), let $A=A_1\cup A_2$ be a union of
two full sieves, with $A_0=A_1\cap A_2$, and let
$B=B_1\cup B_2$ be another such union. Let $f:A\to B$ carry
each $A_i$ into $B_i$. Every simplex of $NA$ lies in one of
$NA_1,NA_2$: its last vertex lies in one of the sieves, and all
its sources then lie there. Consequently
$NA=NA_1\cup_{NA_0}NA_2$, a homotopy pushout of simplicial
sets along monomorphisms. The same holds for $B$. Applying
$\Fun(-,\mathcal C)$ to their groupoid completions gives
pullback squares of infinity-categories. If the three maps for
$i=0,1,2$ are equivalences, so is the fourth. This is precisely
Loc~(6a); it is not either reverse-excision assertion.

For a filtered diagram of arrows in $\Cat$, the nerve commutes
with its filtered colimit, since a finite string of composable
arrows and every equality between such strings occur at some
stage. Filtered colimits of simplicial sets compute the filtered
colimit of the represented spaces. Thus the analogue of (B.2.3)
identifies restriction at the colimit with a limit of restrictions.
Limits preserve equivalences and fully faithful functors. A.8
therefore gives Loc~(8) for $\Wlc$ and for $\Wone$ in the
enhanced case. For the precise component condition on $\Wone$,
G.2 and $\Wone\subseteq\Wc$ prove sufficiency of nonthinness.
If $\D(e)$ is thin up to equivalence, the fold $e\amalg e\to e$
belongs to $\Wc$, and hence to $\Wone$ by A.3, but does not
belong to $\W_0$. This also covers the terminal coefficient
category. Thus nontriviality suffices for Loc~(5) for $\Wlc$,
whereas the strictly stronger nonthinness condition is exact
for $\Wone$ and $\Wc$.

For $\Wc$, B.1 and
\citet[Corollary~4.2.22]{Cisinski2006} give filtered-colimit
closure. For direct excision, take the fundamental sublocalizer
generated by the three given arrows. It is accessible by
definition, is contained in $\Wc$, and corresponds to the
simplicial model structure with monomorphisms as cofibrations
\citep[Theorems~1.4.3 and~4.2.15]{Cisinski2006}. Its usual
cofibrant gluing lemma applies to the two nerve pushouts above.
This proves Loc~(6a) without an accessibility assumption on
$\Wc$ itself. Its sums and component
exceptions are the explicit represented-lattice example G.2.
By B.3, $\Wlc$ is also a fundamental localizer for an arbitrary
derivator. The same generated-sublocalizer gluing argument and
Corollary~4.2.22 give Loc~(6a) and Loc~(8) without enhancement.
For $\Wone$, use instead the fully faithful version of
Grothendieck-construction descent proved in B.3. A natural
transformation of diagrams $A_i\to B_i$ whose components are
in $\Wone$ induces a map of their Grothendieck opfibrations
with fully faithful restriction on locally constant objects.
For a filtered diagram, the comparison from its Grothendieck
construction to its ordinary colimit is in $\Winfty$: the
homotopy-colimit formula (B.2.2) and preservation of filtered
colimits by the nerve identify this comparison with the
usual homotopy-colimit-to-colimit equivalence of simplicial
sets. For a diagram $A_1\leftarrow A_0\to A_2$ of the
sieves above, the same comparison is in $\Winfty$ because
the nerve union is a pushout along monomorphisms. In both
cases these comparisons induce equivalences on local systems
by A.2. Conjugating the fully faithful Grothendieck-construction
restriction by them proves respectively Loc~(8) and Loc~(6a)
for $\Wone$, without using two-out-of-three for that class.
\end{proof}

\subsection{Reverse excision and the three meanings of fibration}

\begin{claim}{B.6}{Refuted}{The reverse Mayer--Vietoris axioms}
For a derivator $\D$ with $\D(e)\not\simeq e$, the class
$\Wlc$ satisfies neither Loc~(6b) nor Loc~(6c) of Chapter XVI.
If $\D(e)\simeq e$, all arrows belong to $\Wlc$ and both
axioms hold. For $\Wone$ of a nontrivial derivator,
Loc~(6b) always fails, whereas Loc~(6c) holds exactly when
$\D(e)$ is a preorder up to equivalence.
Thus the unqualified Loc~(6), which includes
Loc~(6b), fails for every nontrivial coefficient theory, even
for spaces and ordinary weak homotopy equivalences. In contrast,
each constant homology localizer $W_E$ satisfies all three
Mayer--Vietoris implications Loc~(6a,b,c).
\end{claim}
\begin{proof}
Write a Mayer--Vietoris square in the order
$(X_0,X_1,X_2,X_3)$, where $X_1,X_2$ are full sieves of
$X_3$, their union is $X_3$, and $X_0=X_1\cap X_2$.
For a map of two such squares, the manuscript's precise
implications are
\[
\begin{array}{ll}
 (6b):& f_1,f_2,f_3\in W\Longrightarrow f_0\in W,\\
 (6c):& f_0,f_3\in W\Longrightarrow f_1,f_2\in W.
\end{array}
\tag{B.6.1}
\]
We give separate countercubes, so no implication from (6c)
to (6b) is being assumed.

For any small category $A$, let $C_v=A\star\{v\}$ and
$C_w=A\star\{w\}$ be its two terminal cones and put
$\Sigma A=A\star\{v,w\}$, with no arrows between $v$ and
$w$. Its objects are those of $A$ and $v,w$. Its arrows
are those of $A$, the two new identities, and exactly one
arrow $a\to v$ and one arrow $a\to w$ for every $a\in A$.
Composing any arrow of $A$ with either terminal arrow gives
the terminal arrow from its source; these and the compositions
in $A$ specify the entire composition law. The inclusions
$C_v,C_w\subset\Sigma A$ are full sieves with intersection
$A$. Their nerves form a homotopy pushout, and
$B\Sigma A\simeq\Sigma BA$ is the unreduced suspension.
It suffices to find
\[
 A\longrightarrow e\notin\Wlc,
 \qquad \Sigma A\longrightarrow e\in\Wlc.
\tag{B.6.2}
\]

First suppose that $\D(e)$ has distinct parallel morphisms
$\alpha,\beta:M\to N$. Choose a nontrivial integrally acyclic
group $G$ as constructed in B.10, and take $A=BG$, the
one-object category with arrows $G$ and group multiplication
as composition. The CW space $BG$ is connected and acyclic.
Its suspension is simply connected and acyclic, and therefore
weakly contractible by homological Whitehead. Thus the second
condition of (B.6.2) holds by A.2.

For the first condition, let $a:e\to BG$ be its object.
The coherent object $a_!M\in\D(BG)$ has value
$\coprod_{g\in G}M$ at that object; its $G$-action is regular
permutation of the summands. This is a pointwise Kan-extension
calculation, and constructs a coherent object rather than
merely postulating a diagram in $\D(e)$. A nonidentity
permutation is detected by a map to $N$ assigning $\alpha$
to one summand and $\beta$ to all the others. Consequently
its action is nontrivial. It cannot be isomorphic to a
constant object, whose action is trivial. Every object of
$\D(BG)$ is locally constant since every arrow of $BG$ is
invertible. Thus $\D(e)\to\LC_\D(BG)$ is not essentially
surjective, proving the first condition.

It remains to treat a nontrivial $\D(e)$ which is a preorder
up to equivalence. For every nonempty group $G$, a nonempty
coproduct of copies of an object $M$ of this category is
isomorphic to $M$. The counit
$a_!a^*F\to F$, for $F\in\D(BG)$, is therefore an
isomorphism on its sole point evaluation, and hence an
isomorphism by Der~2. Likewise the canonical map
$a_!M\to p^*M$, where $p:BG\to e$, is an isomorphism.
The adjunction for $a$ now gives
\[
 \Hom_{\D(BG)}(p^*M,p^*N)
 \simeq\Hom_{\D(BG)}(a_!M,p^*N)
 \simeq\Hom_{\D(e)}(M,N).
\]
Thus every $BG\to e$ belongs to $\Wlc$. Take instead
$A=e\amalg e$. Its projection is not in $\Wlc$: the
diagonal $\D(e)\to\D(e)^2$ is not essentially surjective
because $\D(e)$ has two nonisomorphic objects. The category
$\Sigma A$ is the four-object poset
$a,b<v,w$. Its nerve is a circle, hence is weakly equivalent
to $B\Z$, and the preceding argument proves the second
condition of (B.6.2). Explicitly, a weak equivalence
$\Sigma A\to B\Z$ sends $a\to v$ to $1\in\Z$, the
other three nonidentity arrows to $0$, and every object to
the unique object; its map on the unique cycle has degree one.

For Loc~(6b), map the square
$(A,C_v,C_w,\Sigma A)$ to the identical double-cone
construction on $e$. Each cone has a terminal object, and
both total categories are $\Wlc$-aspherical by (B.6.2).
Thus $f_1,f_2,f_3$ belong to $\Wlc$, whereas
$f_0:A\to e$ does not. For Loc~(6c), use the map
\[
 (A,A,C_w,C_w)\longrightarrow(A,C_v,C_w,\Sigma A).
\tag{B.6.3}
\]
Both quadruples are Mayer--Vietoris squares, allowing one
sieve to equal the whole category. Here $f_0=\id_A$,
$f_2=\id_{C_w}$, and $f_3:C_w\to\Sigma A$ belong to
$\Wlc$, but $f_1:A\to C_v$ does not, by two-out-of-three
with $C_v\to e$. This directly disproves (6c).
The finite thin-coefficient instance, including all objects,
arrows, compositions, functoriality and sieve checks, is
verified by \texttt{code/B\_verify\_reverse\_excision.py};
its complete output is included.

For the asserted $1$-complete comparison, first assume
that $\D(e)$ is nonthin. Reverse the intersection map in
the first cube: map the double cone of $e$ to the double
cone of the acyclic $BG$. The three cone and total maps
are weak homotopy equivalences. But restriction along
$e\to BG$ is not full: with the objects $M,N$ chosen
above, it induces the nonbijective diagonal
\[
 \Hom(M,N)\longrightarrow\Hom(M,N)^G
\]
on morphisms from the free coherent $G$-object $a_!M$
to the constant object $N$. This refutes (6b).

For (6c), choose a nontrivial acyclic group $H$ and,
by B.10, an acyclic simple group $G$ containing $H$
properly. Its larger cardinality guarantees properness
of the inclusion. Let $Y_1$ be the categorical mapping
cylinder of $BH\to BG$. It has two objects $h,g$,
endomorphism groups $H,G$, and
$\Hom(h,g)=G$, with $\Hom(g,h)=\varnothing$.
Composition on the cross-arrows is left multiplication
by $G$ and right multiplication by the specified image
of $H$. This specifies all its arrows and compositions.
The full subcategory $BH$ is a sieve; projection
$Y_1\to BG$ is a weak homotopy equivalence, as seen
from its terminal-index Grothendieck construction.
Let $C_H=BH\star\{w\}$ and glue it to $Y_1$ along
$BH$, giving $Y_3$. Besides the arrows just listed,
$Y_3$ has $\id_w$ and a unique $h\to w$, fixed by
precomposition with $H$, and no arrows $g\to w$ or
out of $w$. Its two pieces $Y_1,C_H$ are full sieves.
Its classifying space is the homotopy cofiber of
$BH\to BG$. Van Kampen gives
$\pi_1BY_3=G/\langle\!\langle H\rangle\!\rangle=1$
by simplicity; Mayer--Vietoris gives zero reduced
integral homology since both groups are acyclic.
Thus $BY_3$ is weakly contractible.

Use the map of squares
\[
 (BH,BH,C_H,C_H)\longrightarrow(BH,Y_1,C_H,Y_3).
\tag{B.6.5}
\]
Its $f_0$ is an identity and its $f_3$ is a weak
homotopy equivalence. Its $f_1$ is not $1$-complete.
Indeed, restriction of the free coherent $G$-object
to $BH$ is a coproduct of free coherent $H$-objects,
one for each coset in $G/H$. To see this at the
coherent level, choose coset representatives, use
the corresponding arrows of $BG$ to construct the
maps from the free $H$-objects by Kan adjunction,
and apply point evaluation: their disjoint regular
orbits partition $G$, so that map is an isomorphism
by Der~2. The morphism comparison to the constant
object $N$ is consequently
\[
 \Hom(M,N)\longrightarrow\Hom(M,N)^{G/H},
\]
again a nonbijective diagonal. Since $Y_1\to BG$
is a weak equivalence, the same failure holds for
$BH\to Y_1$. This proves failure of (6c) for
$\Wone$ with nonthin coefficients, including all
the nonzero stable coefficient categories.

There is a different answer for thin coefficients.
If $P=\D(e)$ is a nontrivial preorder up to equivalence,
then every $\D(A)$ is thin as well: its diagonal
$F\to F\times F$ is pointwise an isomorphism,
hence an isomorphism by Der~2, which forces every
Hom set to have at most one element. Let
$u:\operatorname{Ob}(A)_{\mathrm{disc}}\to A$ be
the inclusion of its discrete object set. The
unit $F\to u_*u^*F$ is a pointwise isomorphism.
At $a$ the right-hand side is the product of the
values $F_b$ over arrows $b\to a$, in presheaf
variance. The value $F_a$ maps to every factor by
the underlying diagram, and occurs as the identity
factor; in a preorder these two facts make the
unit an isomorphism. The Kan adjunction therefore
makes the underlying-diagram functor fully faithful.
Conversely, for a presheaf of values in $P$, applying
$u_*$ to its values has exactly those point values
by the same product calculation. Thus
$\D(A)\simeq\Fun(A^{\op},P)$, and
\[
 \LC_\D(A)\simeq P^{\pi_0A},\qquad
 \Wone=\{f:\pi_0f\text{ is surjective}\}.
\tag{B.6.6}
\]
Here the last equivalence uses nontriviality of the
complete preorder $P$: an omitted target component
is detected by its distinct initial and terminal
objects. Map the double cone of the empty category
to the double cone of $e$. Its intersection map
$\varnothing\to e$ is not in $\Wone$; both cone
maps are weak equivalences, and the total map from
two points to a connected category is surjective
on components. This refutes (6b).

To verify (6c) in this case, suppose $f_0,f_3$
are surjective on components and let $C$ be a
component of $Y_1$. If $C$ meets $Y_0$, surjectivity
of $\pi_0f_0$ gives an object of $X_0\subset X_1$
mapping into $C$. If it does not, the sieve-cover
description of the nerve shows that $C$ is a whole
component of $Y_3$ disjoint from $Y_2$.
Surjectivity of $\pi_0f_3$ then gives an object
of $X_3$ mapping into $C$; it must lie in $X_1$,
since an object of $X_2$ maps into $Y_2$.
Thus $\pi_0f_1$ is surjective, and the symmetric
argument proves the same for $f_2$. This proves
the exact thin-coefficient exception for (6c).

If $\D(e)\simeq e$, G.1 shows $\D(A)\simeq e$ for
every $A$, and the exception in the statement follows.
Finally, a sieve cover gives a homotopy-pushout square of
nerves. Applying $E\wedge\Sigma^\infty_+$ gives the exact
triangle
\[
 E\wedge\Sigma^\infty_+BX_0\longrightarrow
 (E\wedge\Sigma^\infty_+BX_1)\oplus
 (E\wedge\Sigma^\infty_+BX_2)\longrightarrow
 E\wedge\Sigma^\infty_+BX_3.
\tag{B.6.4}
\]
In a map of exact triangles, equivalence of any two terms
implies equivalence of the third. An equivalence of the
displayed block-diagonal middle map implies equivalence of
each block, since each is a retract. This proves (6a), (6b)
and (6c) for $W_E$, in particular for the constant class
of $D(k)$. This stable excision argument does not apply to
the $2$-complete class of $D(k)$.
\end{proof}

\begin{claim}{B.7}{Proved}{Classification of the fibration axioms}
Let $W$ be a fundamental localizer with
$W\subseteq W_0$. Define $\mathcal F_W$ to be the class
of Grothendieck fibrations whose transport functors belong
to $W$, and define $\operatorname{Fib}_W$ by preservation
of all $W$-equivalences under strict categorical pullback.
Then each of Loc~(7a), Loc~(7b) and Loc~(7c) holds if and
only if $W=W_\infty$. In particular this applies to
$W=\Wlc$ for every nontrivial derivator, by B.3 and G.1.
For the trivial derivator all three axioms hold.

If $W$ is accessible, Loc~(7F), with the manuscript's
class $\operatorname{Fib}_W$, holds if and only if $W$
is proper, equivalently if and only if it is generated
as a fundamental localizer by a set of projections
$A\to e$. For a nonzero associative unital ring $k$ this
specializes to the complete classification
\[
\begin{split}
 \operatorname{Loc}(7F)\text{ for }W_{D(k)}^{\mathrm{lc}}
 &\Longleftrightarrow
  \text{no rational prime is invertible in }k\\
 &\Longleftrightarrow W_{D(k)}^{\mathrm{lc}}=W_\infty.
\end{split}
\tag{B.7.1}
\]
The constant $k$-homology localizer satisfies none of
Loc~(7a,b,c,F), for every nonzero $k$.
For $\Wone$, which need not be a localizer,
Loc~(7c) nevertheless holds for every derivator.
With the enhancement of C.2, Loc~(7a) and the
factorization assertion of Loc~(7F) hold as well;
Loc~(7b) fails for sets, groupoids, every finite
$r$-type theory, every nonzero $D(k)$, every nonzero
$L_E\Sp$, and every nontrivial thin coefficient
theory. These facts do not make $\Wone$ the
weak-equivalence class of a Brown category.
\end{claim}
\begin{proof}
\emph{The exact axioms.}
Chapter XVI, ms.~38--39, states
\[
\begin{array}{ll}
 (7a):&\mathcal F_W\subseteq\operatorname{Fib}_W;\\
 (7b):& f:X\to Y\text{ over }S,\quad
 X\to S,Y\to S\in\mathcal F_W,\quad f\in W\\
 &\hspace{1.5em}\Longrightarrow
    T\times_Sf\in W\text{ for every }T\to S;\\
 (7c):&\mathcal F_W\cap W\subseteq W^{\mathrm u};\\
 (7F):&\text{every functor factors as }p i,
       \quad i\in W,\quad p\in\operatorname{Fib}_W.
\end{array}
\tag{B.7.2}
\]
The word ``cofibrations'' in the manuscript's formulation
of (7F) is a typographical slip: the surrounding definition
is the category-of-fibrant-objects factorization axiom.
The class called \emph{weak $W$-fibrations} in
\citet[Definition~6.1.13]{Cisinski2006} is precisely
$\operatorname{Fib}_W$ here. It is not the manuscript's
weaker class obtained by testing only contractible bases.

\emph{A pullback lemma.}
Suppose that factorizations in (7F) exist. If
$f:X\to Y$ is in $W$ over $S$ and both structural maps
belong to $\operatorname{Fib}_W$, then every pullback
of $f$ over $S$ is in $W$. Indeed, factor $t:T\to S$
as $T\xrightarrow{a}Z\xrightarrow{b}S$, with
$a\in W$ and $b\in\operatorname{Fib}_W$.
The map $f_Z$ is in $W$ by the definition of $b$.
The maps $X_T\to X_Z$ and $Y_T\to Y_Z$ are in $W$
by the definition of the pullbacks of the two structural
maps. Two-out-of-three in the resulting square gives
$f_T\in W$. In particular
\[
 \operatorname{Fib}_W\cap W\subseteq W^{\mathrm u}.
\tag{B.7.3}
\]

\emph{Why (7a) implies (7b).}
Every functor $u:A\to S$ factors as a weak homotopy
equivalence followed by a member of $\mathcal F_{W_\infty}$.
For completeness, factor it first as
$A\to P\xrightarrow{p}S$ in the ordinary Thomason model
structure. Its model fibration $p$ is universally
$W_\infty$-preserving, by right properness. Form the
Grothendieck fibration $q:\Phi(P)\to S$ of the
contravariant diagram $s\mapsto s\backslash P$.
An object is $(s,x,\beta:s\to p(x))$; an arrow to
$(s',x',\beta')$ is a pair $(a:s\to s',h:x\to x')$
with $p(h)\beta=\beta'a$. The inclusion
$j(x)=(p(x),x,\id)$ is a weak homotopy equivalence:
forgetting $s,\beta$ gives a retraction $\epsilon$,
and the arrows $(\beta,\id)$ give
$\id\Rightarrow j\epsilon$.
For $s\to s'$, the transition
$s'\backslash P\to s\backslash P$ is the pullback
by $p$ of $s'\backslash S\to s\backslash S$.
The latter is a weak equivalence between categories
with initial objects, so the transition is a weak
equivalence. This proves the asserted factorization.

Since $\mathcal F_{W_\infty}\subseteq\mathcal F_W$,
(7a) supplies (7F), using these factorizations. Applying
the pullback lemma then gives (7b). Taking $Y=S$ and
$q=\id_S$ gives (7b)$\Rightarrow$(7c).

\emph{Why (7c) forces the minimal localizer.}
The exact hypotheses of
\citet[Theorem~6.4.22]{Cisinski2006} are:
(a) $W\subseteq W_0$; and (b) a $W$-locally constant
functor is a $W$-equivalence if and only if it is
$W$-aspherical. No accessibility hypothesis occurs.
Here ``$W$-locally constant'' means that the transitions
$A/s\to A/s'$ are in $W$.
Given such a functor $u:A\to S$, replace it by the
Grothendieck opfibration
\[
 q:\int_{s\in S}(A/s)\longrightarrow S.
\]
The projection to $A$ is a weak homotopy equivalence,
with section $a\mapsto(u(a),a,\id)$ and the evident
natural comparison; that comparison also gives
$q\in W\Longleftrightarrow u\in W$.
Its fibres are $A/s$ and its transports belong to $W$.
Apply (7c) to its opposite, a Grothendieck fibration.
Fundamental localizers are invariant under opposites.
We conclude that every $A/s$ is $W$-aspherical.
Thus a $W$-locally constant $W$-equivalence is
$W$-aspherical. The converse is Theorem A.
The cited theorem proves $W=W_\infty$.

For $W_\infty$, a Grothendieck fibration with weakly
equivalent transports computes its strict base changes
as homotopy pullbacks, by Quillen's Theorem B and its
base-change criterion
\citep[\S1, Theorem B and its corollary]{Quillen1973};
see also \citet[Proposition~2.7]{Rezk1998Sharp}.
It therefore preserves weak equivalences under every
base change, giving (7a), and hence (7b,c). This completes
their classification. The same result holds for every
nontrivial fundamental localizer, without the explicit
condition $W\subseteq W_0$, by
\citet[Remark~6.4.23]{Cisinski2006}. The only additional
nontrivial case is the gross localizer; the inclusion of
one point into two discrete points refutes (7c) there.

\emph{Why (7F) is exactly properness.}
Let $W_{\mathrm{asph}}$ denote the fundamental localizer
generated by all projections $A\to e$ which belong
to $W$. It is contained in $W$ and has exactly the same
aspherical categories and aspherical functors
\citep[\S6.1.7]{Cisinski2006}.
Suppose (7F). For any $f:X\to Y$ in $W$, factor its
graph as
\[
 X\xrightarrow{i}P\xrightarrow{q}X\times Y,
 \qquad i\in W,\quad q\in\operatorname{Fib}_W.
\]
The projections from a product belong to
$\operatorname{Fib}_W$, since a fundamental localizer
is stable under products. Consequently
$r=\operatorname{pr}_Xq$ and $p=\operatorname{pr}_Yq$
belong to $\operatorname{Fib}_W$. Since $ri=\id_X$
and $pi=f$, both $r,p$ are also in $W$.
By (B.7.3) they are universally in $W$.
For each $y\in Y$, pull $p$ back to $Y/y$. The result
$P/y\to Y/y$ is in $W$, and its target is contractible.
Thus $P/y$ is $W$-aspherical, equivalently
$W_{\mathrm{asph}}$-aspherical. Theorem A for
$W_{\mathrm{asph}}$ gives $p\in W_{\mathrm{asph}}$.
The same argument gives $r\in W_{\mathrm{asph}}$.
Then $ri=\id_X$ gives $i\in W_{\mathrm{asph}}$, and
$f=pi\in W_{\mathrm{asph}}$. We have proved
\[
 \operatorname{Loc}(7F)\Longrightarrow W=W_{\mathrm{asph}},
\tag{B.7.4}
\]
without an accessibility assumption.

If $W$ is accessible, Proposition~6.1.8 and
Theorem~6.1.11 of \citet{Cisinski2006} identify this
equality with properness, or equivalently generation
by a set of aspherical projections. Conversely, for
proper $W$ the Thomason model structure is right proper
\citep[Theorem~5.2.15]{Cisinski2006}. Its model
fibrations belong to $\operatorname{Fib}_W$
\citep[Example~6.1.14]{Cisinski2006}, and its
factorization axiom supplies (7F).
This proves the claimed equivalence.
Here properness is right properness of the associated
simplicial model structure; the cited theorem gives the
equivalent Thomason-model condition. Literal Loc~(9),
accessibility of the full subcategory of arrows, must not
be confused with the additional assertion
$W=W_{\mathrm{asph}}$.

\emph{An obstruction to properness.}
We use the precise statement of
\citet[Lemma~9.4.5]{Cisinski2006}: if $W\subseteq W_n$
and $W$ is proper, every simplicial $W$-equivalence
whose target is $W$-fibrant is a \emph{strong}
$n$-equivalence. By Scholie~9.2.18 there, this includes
surjectivity on $\pi_{n+1}$, in addition to isomorphisms
through degree $n$. It is enough below to exhibit a
$W$-equivalence between connected spaces, with
$W$-local target and nonsurjective $\pi_{n+1}$.
A Kan representative of a $W$-local space is
$W$-fibrant: restriction against a $W$-equivalence
which is a cofibration is a weak equivalence and a
Kan fibration of mapping spaces, hence a trivial
fibration. This verifies the target hypothesis in
each application, rather than omitting it.

\emph{Characteristic zero rings.}
Let $S$ be the set of rational primes invertible in
$k$, and set $R=\Z[S^{-1}]$. By the arithmetic test
of C.9, $k$-homology equivalences are exactly
$R$-homology equivalences: both are detected by
$\Q$ and by $\F_p$ for $p\notin S$.
If $S$ is empty, C.9 gives
$W_{D(k)}^{\mathrm{lc}}=W_\infty$, which is proper.
If $S$ is nonempty, consider
\[
 j:K(\Z,2)\longrightarrow K(R,2).
\tag{B.7.5}
\]
Its homotopy fibre is $K(R/\Z,1)$. The group $R/\Z$
is the filtered union of finite groups whose orders
are invertible in $R$. Averaging makes their positive
$R$-homology zero, and the bar construction commutes
with filtered unions. Thus this fibre is $R$-acyclic.
The Serre homology spectral sequence proves that $j$
is an $R$-homology equivalence. Both its spaces are
simply connected, so C.7 makes it a
$W_{D(k)}^{\mathrm{lc}}$-equivalence.

The target is local for $R$-homology equivalences.
Indeed, the homotopy groups of
$\Map(X,K(R,2))$ are $H^{2-i}(X;R)$ for $0\leq i\leq2$;
these are invariant under $R$-homology equivalences,
by applying derived $R$-linear Hom into $R$ to
the singular-chain equivalence. It is consequently
local for the smaller class $W_{D(k)}^{\mathrm{lc}}$.
But $\pi_2(j):\Z\to R$ is not onto. Since C.7
also gives $W_{D(k)}^{\mathrm{lc}}\subseteq W_1$,
the preceding obstruction with $n=1$ proves that
this localizer is not proper.

\emph{Positive characteristic rings.}
Let $d=\operatorname{char}(k)>0$, and let $P$ be
the nonempty finite set of primes dividing $d$.
The nonunits among rational primes are precisely
the members of $P$. C.9 identifies $k$-homology
equivalences with simultaneous $\F_p$-homology
equivalences for $p\in P$. Put
\[
 M=\prod_{p\in P}\Z_p,
 \qquad j:K(\Z,2)\longrightarrow K(M,2),
\tag{B.7.6}
\]
where $\Z\to M$ is diagonal and $\Z_p$ denotes
the additive group of $p$-adic integers, regarded here as a discrete group.
For each $p\in P$, multiplication by $p$ is an
automorphism of $A=M/\Z$. For surjectivity, given
$m\in M$, choose an integer $z$ congruent to its
$p$-component modulo $p$. Then $(m-z)/p\in M$:
division by $p$ is allowed in the other components.
For injectivity, if $pm$ is the diagonal integer
$z$, integrality in the $p$-component implies
$p\mid z$, and then $m$ is itself diagonal.

For clarity, an abelian group uniquely divisible
by $p$ is $\F_p$-acyclic as a group. Such a group
has no $p$-torsion. Each finitely generated subgroup
$H$ is therefore a sum of a free abelian group and
a finite group of order prime to $p$, so
$H_*(BH;\F_p)$ is the exterior algebra on its first
homology. Enlarge $H$ to the finitely generated
subgroup $p^{-1}H$. Its inclusion induces zero on
$H_1(-;\F_p)$, since under the isomorphism
$p^{-1}H\simeq H$ it is multiplication by $p$.
It therefore induces zero in every positive
homology degree. Filtered colimits of the bar
complexes prove the assertion. Hence $K(A,1)$,
the homotopy fibre of (B.7.6), is $\F_p$-acyclic
for every $p\in P$. The same Serre argument proves
that $j$ is in $W_{D(k)}^{\mathrm{lc}}$.

The target in (B.7.6) is local for $k$-homology
equivalences. To verify this rather than assume
it, an $\F_p$-homology equivalence induces all
$\F_p$-cohomology isomorphisms by vector-space
duality. The coefficient short exact sequences
\[
 0\longrightarrow\Z/p^{r-1}
 \longrightarrow\Z/p^r\longrightarrow\F_p
 \longrightarrow0
\]
and their long exact cohomology sequences imply
the same assertion for $\Z/p^r$. Thus
$K(\Z/p^r,2)$ is $\F_p$-homology local.
Local spaces are closed under homotopy limits,
and the surjective inverse system of these
coefficient groups gives
\[
 \holim_rK(\Z/p^r,2)\simeq K(\Z_p,2);
\]
the possible derived-limit group is zero by
surjectivity. Taking the finite product over $P$
proves the required locality for the simultaneous
homology class, and hence for
$W_{D(k)}^{\mathrm{lc}}$. The map
$\pi_2(j):\Z\to M$ is not surjective, since $M$
is uncountable. The $n=1$ obstruction again
proves nonproperness. Together with B.4 and
(B.7.4), this proves (B.7.1) for every nonzero ring.

\emph{The constant homology class.}
For every nonzero $k$, it contains $BG\to e$ for
the nontrivial acyclic group used in B.6, so it
strictly contains $W_\infty$. It preserves
components, and the already proved classification
therefore refutes (7a,b,c). It is also never proper.
In characteristic zero with $S$ empty it equals
the integral homology localizer, whose
nonproperness is
\citet[Scholie~9.4.7]{Cisinski2006}.
When $S$ is nonempty, use
$K(\Z,1)\to K(R,1)$ and the obstruction with
$n=0$. Its target is $R$-homology local by the
same cohomology argument. Its $R$-homology is an
isomorphism: $R$ is the filtered union of its
cyclic subgroups, the bar homology commutes with
that union, and the only positive group homology
of $\Z$ is in degree one. The map on $\pi_1$ is
not surjective.

In positive characteristic use
$K(\Z,1)\to K(M,1)$. The preceding locality proof
works in degree one as well. The torsion-free
abelian group $M$ is a filtered union of finitely
generated free abelian groups, whence
\[
 H_*(BM;\F_p)\simeq
 \bigwedge_{\F_p}^*(M/pM)
 \simeq\bigwedge_{\F_p}^*\F_p
 \quad(p\in P).
\]
The diagonal $\Z\to M$ induces the displayed
isomorphism in degree one, and hence is a
$k$-homology equivalence. Again it is not onto
on $\pi_1$. Lemma~9.4.5 with $n=0$ proves
nonproperness. Accessibility of these homology
localizers is B.4, or
\citet[Theorem~9.4.2 and Corollary~9.4.3]{Cisinski2006}.
Consequently (7F) fails in every nonzero-ring case.

\emph{The intermediate class $\Wone$.}
The fundamental-localizer implications above must
not be applied to this class, because it need not
satisfy two-out-of-three. We instead check each
assertion directly. For variance, put
$\D^{\op}(A)=\D(A^{\op})^{\op}$. Then
$f\in W^1_{\D}$ if and only if
$f^{\op}\in W^1_{\D^{\op}}$.
It therefore suffices to prove the assertion for
Grothendieck opfibrations and apply that proof
to $\D^{\op}$ when the manuscript uses
Grothendieck fibrations. This dualizes both the
shape and coefficient derivator; it does not
assume shape-opposite invariance for a fixed
bare derivator. For opfibrations, right Kan
extensions in the presheaf convention have
the strict-fibre formula of B.3.
If $p:E\to S$ has transports in $\Wone$, then
$p_*$ sends locally constant objects to locally
constant objects, and does so after every base
change. Indeed, each transport is in the
cohomological-equivalence class by A.3; the
comparison of the two fibre cohomology objects
is an isomorphism for every locally constant
coefficient. Pointwise Kan extension gives
exactly this comparison on the base arrow.
Consequently $p^*$ and $p_*$ restrict to an
adjunction on locally constant categories.

If also $p\in\Wone$, the unit of this adjunction
is an isomorphism. Apply it to the constant
object of value $M\in\D(e)$ and evaluate at
$s\in S$. The strict-fibre formula gives
\[
 M\xrightarrow{\sim}
 (p_{E_s})_*p_{E_s}^*M.
\tag{B.7.8}
\]
Thus $E_s\to e$ is a constant-coefficient
equivalence, equivalently a member of $\Wone$
for this point target, by A.3. Every base change
of $p$ is again an opfibration with the same
fibres. The fully faithful case of B.3 proves
that each such base change belongs to $\Wone$.
This is literal Loc~(7c), without any two-out-of-three
argument.

Under the enhancement, C.2 identifies
$\Wone$ with the cohomological class. The
all-local-coefficients direct-image property
just proved and its Beck--Chevalley formula
make $p$ universally preserve that class, by
the adjunction argument of D.2. Explicitly,
for $f:T\to S$ in the cohomological class and
$F\in\LC(E)$, cohomology on $E$ equals
cohomology on $S$ with coefficient $p_*F$.
This coefficient is locally constant, so $f$
preserves it; base change identifies the result
with the comparison on $T\times_SE$.
This proves Loc~(7a) for $\Wone$ in the
enhanced case. For the factorization assertion,
factor any functor in the ordinary Thomason
model structure. Its first map is in
$W_\infty\subseteq\Wone$. Its second map
is an ordinary sharp fibration, whose strict
base changes compute homotopy pullbacks. C.2
makes $\Wone$ stable under homotopy pullback,
so this second map is in
$\operatorname{Fib}_{\Wone}$. This proves
the asserted factorization, without claiming
the Brown axioms that require two-out-of-three.

Finally choose a pointed category $Z$ with
$Z\to e\in\Wone$ and $e\to Z\notin\Wone$.
Let $R=Z\star\{t\}$ and $p:R\to[1]$ be
the Grothendieck opfibration of $Z\to e$.
The chosen point gives a section
$s:[1]\to R$. Both structural maps $p$ and
$\id_{[1]}$ have transports in $\Wone$.
The section is a weak homotopy equivalence
because both total categories have terminal
objects, but its fibre at $0$ is the excluded
map $e\to Z$. Taking opposites gives the
literal Grothendieck-fibration version. This
refutes (7b), not (7c): the structural map $p$
itself is universally in $\Wone$ by (B.7.8).
For $r$-types take a category model of
$K(A,r+1)$ with $A\ne0$; for groupoids use
$r=1$. For nonzero stable coefficients take
$Z=BG$ for the nontrivial integrally acyclic
group of B.10: its suspension is contractible,
so constants satisfy the fibre criterion C.2,
whereas $e\to BG$ fails full faithfulness by
the regular-action argument in B.6. For a
nontrivial thin coefficient category take
$Z=e\amalg e$ and use (B.6.6).
All these choices meet the two stated tests.
For spaces $\Wone=W_\infty$ and (7b) holds.
\end{proof}

Here is the resulting comparison for the computed theories.
The zero coefficient theory is understood separately: all
its arrows are equivalences and all listed axioms hold.
\begin{center}
\begin{tabular}{p{0.39\linewidth}ccc}
\hline
Class & Loc (6b,c) & Loc (7a,b,c) & Loc (7F)\\
\hline
$W^{\mathrm{lc}}$ for spaces, $D(\Z)$, spectra
 & no & yes & yes\\
$W^{\mathrm{lc}}$ for sets, groupoids, finite $r$-types
 & no & no & yes\\
$W^{\mathrm{lc}}_{D(k)}$, $k\ne0$, all primes nonunits
 & no & yes & yes\\
$W^{\mathrm{lc}}_{D(k)}$, $k\ne0$, some prime a unit
 & no & no & no\\
Constant $k$-homology, $k\ne0$
 & yes & no & no\\
\hline
\end{tabular}
\end{center}
For sets, groupoids and finite $r$-types the last column
uses C.4 and the properness of every finite truncation
localizer $W_n$, explicitly established in
\citet[\S9.2.1]{Cisinski2006}. Thus (7F) does not imply
(7a,b,c). For a general spectrum $E$, B.4 makes
$W^{\mathrm{lc}}_{L_E\Sp}$ accessible, so the exact
remaining test for (7F) is
\[
 W^{\mathrm{lc}}_{L_E\Sp}
   =(W^{\mathrm{lc}}_{L_E\Sp})_{\mathrm{asph}},
\tag{B.7.7}
\]
equivalently properness. No assertion that (B.7.7)
holds for every $E$ is made; its evaluation beyond
the coefficient theories computed here remains
unresolved in this work. For an arbitrary derivator
without an accessibility theorem, (7F) still implies
aspherical generation by (B.7.4); the converse is
not claimed without that hypothesis.

For an arbitrary unenhanced derivator, two intermediate-class
questions remain unresolved here: whether every Grothendieck
fibration with transports in $\Wone$ belongs to
$\operatorname{Fib}_{\Wone}$ (literal Loc~(7a)), and
whether every functor factors as a member of $\Wone$
followed by a member of $\operatorname{Fib}_{\Wone}$
(the factorization assertion of Loc~(7F)). The first
argument above proves universal preservation of the
cohomological class; identifying its output with
$\Wone$ used the enhancement in C.2. The second
argument used the enhanced homotopy-fibre criterion
to prove homotopy-pullback stability of $\Wone$.
Neither of these two extra steps is supplied just
by the arbitrary-derivator relative Theorem A of B.3.
This precisely limits the two positive enhanced
assertions; it does not leave (7c) unresolved.

Finally, the existence of the model-fibration class
in D.6 is compatible with all the negative entries.
An accessible fundamental localizer always has the
Thomason model structure. Its model fibrations compose,
are stable under strict base change, and detect fibre
equivalences between fibrant objects of a slice.
When the localizer is not proper, they need not belong
to the manuscript's $\operatorname{Fib}_W$.
Replacing that distinction by an identification would
incorrectly turn D.6 into Loc~(7F).


\subsection{A one-object test for truncation localizers}

\begin{claim}{B.8}{Known}{The sphere test, without accessibility}
For every fundamental localizer $\W$ and every integer $n\geq0$,
\[
 \W_n\subseteq\W
 \quad\Longleftrightarrow\quad
 \bigl(\Delta/\partial\Delta^{n+2}\longrightarrow e\bigr)\in\W.
 \tag{B.8.1}
\]
No accessibility hypothesis is needed. In particular, this applies
to $\Wlc$ for every derivator, by B.3.
For $\mathcal C=\mathcal S_{\leq r}$ it holds exactly when
$n\geq r+1$. For $\mathcal C=D(k)$ with $k\ne0$, or for
ordinary spectra, it holds for no finite $n$.
\end{claim}
\begin{proof}
The equivalence is
\citet[Proposition~9.2.17]{Cisinski2006}, which we checked in
the original text. The category of simplices on the right has
classifying space $S^{n+1}$. C.4 identifies the truncated case
with $\W_{r+1}$; inclusion $\W_n\subseteq\W_{r+1}$ is
equivalent to $n\geq r+1$. In the linear cases C.7 and C.10
detect $S^1$ by its fundamental group, and detect every higher
sphere by the nonzero top reduced homology of its universal
cover with the stated coefficients. These are computations of
the specified projection, so they prove the last assertion.
TOTASPH F.28 is confirmed in precisely this form.
\end{proof}

\subsection{Acyclic epimorphisms and cardinal-controlled cones}

\begin{claim}{B.9}{Proved}{An acyclic map is an epimorphism in the unpointed homotopy category}
Let $t:X\to Y$ be a map of small CW spaces whose homotopy
fibres are integrally acyclic, where ``acyclic'' includes
nonempty and connected. Then for every small space $Z$ the
map of unpointed homotopy classes
\[
 [Y,Z]\longrightarrow[X,Z]
 \tag{B.9.1}
\]
induced by $t$ is injective. The homotopy codiagonal
$c:Y\cup_X^hY\to Y$ is a weak homotopy equivalence.
\end{claim}
\begin{proof}
First suppose $Y$ is connected. The homotopy-fibre sequence
makes $X$ connected and $\pi_1(t)$ surjective. Put
$P=Y\cup_X^hY$, realized by the double mapping cylinder, with
its two maps $j_1,j_2:Y\to P$ and its fold $c:P\to Y$.
Van Kampen identifies $\pi_1P$ with
$\pi_1Y*_{\pi_1X}\pi_1Y$. Since $\pi_1(t)$ is onto, the
fold of this amalgam is an isomorphism onto $\pi_1Y$.

An acyclic map induces an isomorphism in homology with every
local system on its target; this is the equivalence of (a) and
(b) in \citet[Definition/Proposition~1.2]{HausmannHusemoller1979}.
For a local system $L$ of abelian groups on $Y$, compare the
Mayer--Vietoris sequence of $P$ with that of the diagram
$Y\leftarrow Y\to Y$, using $t$ at the central term and
identities at the two outer terms. The three component maps
are homology isomorphisms with the pulled-back coefficients.
The induced map $H_*(P;c^*L)\to H_*(Y;L)$ is therefore an
isomorphism. Take $L=\Z[\pi_1Y]$. As $\pi_1c$ is an
isomorphism, this is the integral homology comparison between
the universal covers. Those covers are simply connected CW
spaces. Homological Whitehead, followed by the long exact
sequences of the coverings, makes $c$ a weak equivalence.

Since $cj_1=cj_2=\id_Y$, we have $j_1=j_2$ in the unpointed
homotopy category. If $a,b:Y\to Z$ and $at$ is homotopic to
$bt$, that actual homotopy defines $h:P\to Z$ on the double
mapping cylinder. Consequently $[a]=[h][j_1]=[h][j_2]=[b]$.
For disconnected $Y$, acyclicity gives a bijection on
components, and the argument applies on each component.
The product of these injective maps of sets is injective;
this is a product of unpointed homotopy-class sets of spaces,
not an interchange of an infinite product with components of
ordinary functor categories. The empty case is immediate.
\end{proof}

\begin{claim}{B.10}{Proved}{Acyclic simple overgroups of controlled cardinality}
Let $H$ be a small group and let $\lambda$ be an infinite
small cardinal with $\lambda\geq|H|$. There is an embedding
$H\hookrightarrow A$ into a nonabelian simple integrally
acyclic group with $|A|=\lambda$. Consequently, if
$\lambda>\mu$, every homomorphism $A\to L$ to a group of
cardinality at most $\mu$ is trivial.
\end{claim}
\begin{proof}
We give the two envelope constructions and then alternate them.
For the simple envelope, start with $H*F_\lambda$, so that
the cardinality is exactly $\lambda$. For each pair
$g\ne1,h$ in a current group $L$ of cardinality $\lambda$,
adjoin a separate infinite cyclic free generator $t_{g,h}$
and put, in $L*F(\{t_{g,h}\})$,
\[
 w_{g,h}=g t_{g,h}g^{-1}t_{g,h}^{-1}
          g t_{g,h}^{2}g^{-1}t_{g,h}^{-2}.
 \tag{B.10.1}
\]
Both $w_{g,h}$ and $w_{g,h}h$ have infinite order. The first
is cyclically reduced of length eight in the relevant free
product. For the second, cyclic reduction combines its first
and last $L$-syllables; the result has length eight if $hg\ne1$
and length six if $hg=1$, with every syllable nontrivial.
Thus we may adjoin a stable letter $s_{g,h}$ satisfying
\[
 s_{g,h}^{-1}w_{g,h}s_{g,h}=w_{g,h}h.
 \tag{B.10.2}
\]
The simultaneous HNN embedding theorem
\citep[Theorem~II]{HigmanNeumannNeumann1949} embeds the old
group in this extension, because the identified cyclic
subgroups are infinite cyclic. Equation (B.10.2) puts $h$
in the normal closure of $g$. Iterate this step countably.
In the union any nonidentity $g$ and any $h$ occur together
at some stage; hence the normal closure of $g$ contains $h$.
The union is simple, is nonabelian because it contains
$F_\lambda$, and has cardinality $\lambda$.

For clarity we also verify the needed acyclic envelope. For
a group $L$, start with $L\times L$ and adjoin stable letters
$d,s$ such that, writing $j(h)=(h,1)$ and $k(h)=(1,h)$,
\[
 d^{-1}j(h)d=j(h)k(h),\qquad s^{-1}j(h)s=k(h).
 \tag{B.10.3}
\]
The three copies of $L$ being identified are embedded copies,
so HNN again embeds $L\times L$. Denote the resulting group
by $m(L)$, with embedding $j:L\to m(L)$. This is the
mitosis construction underlying
\citet[Proposition~4.1 and Theorem~4.2]{BaumslagDyerHeller1980}.
Its countable iteration is acyclic, as follows. Over any field
$F$, suppose $f:H_0\to L$ induces zero in positive homology
degrees below $n$. In the K\"unneth decomposition of
$H_n(L\times L;F)$, the diagonal after $f$ has only the two
axis terms: the other terms contain a positive lower-degree
factor killed by $f$. Under (B.10.3), its two axis maps are
both conjugate to $j$, and its diagonal map is also conjugate
to $j$. Hence
\[
 (jf)_*=(jf)_*+(jf)_*\quad\hbox{on }H_n(H_0;F),
\]
so $(jf)_*=0$. Induction shows that the map from any stage
to its $n$-th successor kills positive homology through
degree $n$. Every bar cycle has finite support, so the union
has zero positive homology over every field. Its positive
integral homology groups are torsion by rational coefficients;
the universal coefficient sequence in the next degree with
$\F_p$ coefficients says they have no $p$-torsion, for every
prime $p$. They are therefore zero. This construction preserves
the bound $\max(|L|,\aleph_0)$.

Starting from $H*F_\lambda$, form a simple envelope, then an
acyclic envelope, then a simple envelope, and so on, always
using embeddings. Let $A$ be the union. A nontrivial normal
subgroup of $A$ contains some $g\ne1$; given any $h\in A$,
both lie in a later simple stage, whose simplicity puts $h$
in that normal subgroup. Thus $A$ is simple. Every finite
bar cycle lies in a later acyclic stage and bounds there,
so $A$ is acyclic. Its cardinality is exactly $\lambda$.
Finally the kernel of a homomorphism from a simple group
is trivial or the whole group. An injection into a group
of size $\mu<\lambda$ is impossible, proving the last claim.
\end{proof}

\begin{claim}{B.11}{Proved}{The cardinal-controlled Kan--Thurston statement}
Let $Y$ be a connected $\U$-small CW complex and let $\mu$ be
a cardinal belonging to $\U$. There are a $\U$-small group $G$
and an acyclic map $t:BG\to Y$ such that
$\phi=\pi_1(t):G\twoheadrightarrow\pi_1Y$ is onto, its kernel
$P$ is perfect, and every homomorphism $G\to L$ with
$|L|\leq\mu$ vanishes on $P$. Base points are chosen only to
write $\phi$; acyclicity and the conclusion about the normal
subgroup do not depend on their choice.
\end{claim}
\begin{proof}
Choose a simplicial complex $K$ and a homotopy equivalence
$h:|K|\to Y$; the existence is
\citet[Theorem~2C.5]{C_Hatcher2002}. Choose an infinite
cardinal $\lambda\in\U$ larger than $\mu$ and at least the
number of simplices of $K$. Such a cardinal exists inside the
fixed Grothendieck universe. We construct, simultaneously for
all subcomplexes $L\subseteq K$, a CW space $T(L)$ and a map
$t_L:T(L)\to|L|$, so that inclusions are actual CW inclusions,
every component of $T(L)$ is aspherical, and every inclusion
$T(L')\hookrightarrow T(L)$ is injective on the fundamental
group of each component. All choices attached to a simplex
are made once and used in every subcomplex containing it.

In dimension at most one, put $T(L)=|L|$ and $t_L=\id$.
Graphs are aspherical, and inclusions of subgraphs induce
injections on fundamental groups, as is seen from reduced
edge paths after choosing maximal forests. Suppose the
construction and its inclusion properties have been made
through dimension $n-1$, with $n\geq2$. For each $n$-simplex
$\sigma$ of $K$, the space $T(\partial\sigma)$ is connected
and aspherical. Its fundamental group embeds, by B.10, in
an acyclic nonabelian simple group $C_\sigma$ of cardinality
$\lambda$. Realize this embedding by a cellular map
\[
 a_\sigma:T(\partial\sigma)\longrightarrow BC_\sigma.
 \tag{B.11.1}
\]
The realization is possible because both spaces are connected
$K(\pi,1)$ spaces; one can send their universal covers
equivariantly into the contractible universal cover of
$BC_\sigma$, extending over cells. Let $Q_\sigma$ be its
mapping cylinder, with $T(\partial\sigma)$ at the free end.
Thus $Q_\sigma$ deformation retracts onto $BC_\sigma$,
$T(\partial\sigma)\hookrightarrow Q_\sigma$ is a CW
cofibration, and this inclusion is injective on $\pi_1$.
For a subcomplex $L$ of dimension at most $n$, put
\[
 T(L)=T(L^{(n-1)})\ \cup_{\coprod_{\sigma\in L_n}
                  T(\partial\sigma)}\
                  \coprod_{\sigma\in L_n}Q_\sigma .
 \tag{B.11.2}
\]
On the cylinder of (B.11.1), define $t_L$ by moving
$t_{\partial\sigma}(x)$ linearly inside $\sigma$ to its
barycentre as the cylinder coordinate runs from its free end
to $BC_\sigma$. Send all of $BC_\sigma$ to that barycentre.
This defines an actual continuous map extending the one on
the lower skeleton, and it commutes with all subcomplex
inclusions. All these spaces have at most $\lambda$ cells:
the bar models of the groups do, and finite products and a
countable union preserve this infinite cardinal bound.

We justify both induction assertions, including the
subcomplex injectivity which is needed later. We use the
following elementary graph-of-spaces argument. If connected
aspherical vertex spaces are joined by mapping cylinders
of aspherical edge spaces, and both edge homomorphisms
are injective, the normal form for amalgamated products and
HNN extensions embeds all vertex groups. In the universal
cover the incidence graph is the corresponding tree of
cosets. Vertex and edge inverse images are contractible.
Every finite subtree of this lifted graph of spaces is
contractible by successive attachment of a contractible
vertex space along a contractible edge cylinder. A map
from a sphere has compact image in a finite subcomplex,
and hence lies over a finite subtree. It is nullhomotopic
there. The universal cover is therefore weakly contractible,
and, being a CW space, contractible. This proves
asphericity as well as injectivity; disconnected initial
spaces are handled component by component. It is the
standard argument underlying the amalgamation statement
in \citet[\S1, before Proposition~1.1]{BaumslagDyerHeller1980}.

Apply this first to (B.11.2). Each edge group injects into
the lower-skeleton group by the induction hypothesis and
into $C_\sigma$ by its construction. It follows that
$T(L)$ is aspherical on every component and that its lower
skeleton and all $BC_\sigma$ inject on $\pi_1$.
For the stronger assertion, let $L\subseteq K$ with both
complexes of dimension at most $n$. Form first
\[
 V=T(K^{(n-1)})\cup_{T(L^{(n-1)})}T(L).
 \tag{B.11.3}
\]
This is a graph of spaces with one edge for each component
of $T(L^{(n-1)})$. Its edge injections into
$T(K^{(n-1)})$ are the dimension-$n-1$ induction assertion;
its edge injections into $T(L)$ were proved in the preceding
sentence. The argument just given embeds both vertex groups
in each resulting component of $V$. Next attach the
$Q_\sigma$ for $\sigma\in K_n\setminus L_n$. Each
attaching subgroup injects into $T(K^{(n-1)})$, which already
injects into $V$ and continues to inject at every attachment.
The final space is exactly $T(K)$: both constructions use
the same pushouts and the same chosen $Q_\sigma$. Thus
$T(L)\to T(K)$ is injective on fundamental groups.
This proves the missing subcomplex step without supposing
that an arbitrary map between aspherical spaces is injective.

Take the union over dimensions. Compactness of spheres and
homotopies in CW complexes proves the same assertions for
arbitrary subcomplexes of an infinite-dimensional $K$.
The construction does not introduce new components, so
$T(K)$ is connected when $K$ is. We now check the exact
homological property, not merely constant homology. For
every local system $M$ of abelian groups on $|L|$, induction
on dimension gives
\[
 H_*(T(L);t_L^*M)\xrightarrow{\ \cong\ }H_*(|L|;M).
 \tag{B.11.4}
\]
Indeed, in the relative pair for an $n$-simplex the target
system is constant because $\sigma$ is contractible. The
pair on the source is $(Q_\sigma,T(\partial\sigma))$.
Its long exact homology sequence, acyclicity of
$Q_\sigma\simeq BC_\sigma$, and the induction comparison
for $T(\partial\sigma)\to\partial\sigma$ identify its
relative homology with that of $(\sigma,\partial\sigma)$,
with every constant coefficient group. Excision gives the
direct sum of these relative comparisons for (B.11.2).
The long exact sequence of the skeletal pair proves
(B.11.4). Cellular chains with local coefficients commute
with the filtered union over skeleta, which proves it in
unbounded dimension as well.

Write $G=\pi_1T(K)$ and let $N$ be the normal closure in
$G$ of the embedded groups $C_\sigma$, for all simplices
of dimension at least two. The map $t_K$ kills each of
these groups. Moreover
\[
 G/N\xrightarrow{\ \cong\ }\pi_1|K|.
 \tag{B.11.5}
\]
To verify the specified map, use van Kampen in (B.11.2).
At dimension two, killing $C_\sigma$ kills precisely
the image of $\pi_1T(\partial\sigma)=\pi_1\partial\sigma$,
the usual boundary relator of that two-simplex. At a
dimension $n\geq3$, the inductive version of (B.11.5)
for $\partial\sigma$ says that its group is normally
generated by its lower-dimensional cone groups, because
$\pi_1\partial\sigma=0$. Those groups have already been
killed. No further relation in the quotient is therefore
introduced. This proves (B.11.5) by induction and then by
filtered union. In particular $\pi_1(t_K)$ is onto and
its kernel is exactly $N$. Each $C_\sigma$ is perfect;
its conjugates generate $N$, and each such conjugate lies
in $[N,N]$. Hence $N$ is perfect.

Since $T(K)$ is a connected aspherical CW space, choose
a based homotopy equivalence $b:BG\to T(K)$ inducing the
identity on $G$, and put $t=h\,t_K\,b$. Equation (B.11.4)
and \citet[Definition/Proposition~1.2]{HausmannHusemoller1979}
show that this specified map is acyclic. Its kernel is
$P=N$, transported by the chosen base-point identification.
Every homomorphism $\rho:G\to L$ with $|L|\leq\mu$
restricts trivially to each $C_\sigma$, since that group
is simple and has cardinality $\lambda>\mu$. Thus $\rho$
kills their normal closure $P$. If $K$ has dimension at
most one, there are no cone groups, $P=1$, and the same
conclusion holds. This proves all the stated cases.
\end{proof}

\begin{claim}{B.12}{Proved}{Small ind-objects modulo natural transformations see at most one-types}
Let $\mathcal H$ be the category of $\U$-small categories
whose morphisms are functors modulo zigzags of natural
transformations. Let $I$ be a nonempty $\U$-small filtered
category, and let $K:I\to\mathcal H$ be a diagram, represented
by categories $K_i$. Suppose the presheaf
\[
 H_K(X)=\colim_{i\in I}\Hom_{\mathcal H}(X,K_i)
       =\colim_i\pi_0\Fun(X,K_i)
 \tag{B.12.1}
\]
sends every functor in $\Winfty$ to a bijection. Then it
sends every functor in $\W_1$ to a bijection. The specified
unit of the reflector induces an isomorphism
\[
 (K_i)_{i\in I}\longrightarrow(\Pi_1K_i)_{i\in I}
 \quad\hbox{in }\operatorname{Ind}_{\U}(\mathcal H).
 \tag{B.12.2}
\]
In particular the small-index version of CATHOT E.17 is
settled, with its original indexing category unchanged.
\end{claim}
\begin{proof}
We first identify exactly the homotopy-invariant functor
which (B.12.1) defines. The functor
$q:\mathcal H\to\Hot$ is localization at nerve weak
equivalences: a natural transformation is a functor on
$X\times[1]$, and the two endpoint inclusions become
equal after inverting its projection to $X$. Thus the
localization of $\Cat$ factors through $\mathcal H$
and has the same universal property there.
By hypothesis $H_K$ is the pullback of a presheaf $H'$
on $\Hot$. There is a canonical identification
\[
 H'(Y)\cong\colim_i[Y,BK_i].
 \tag{B.12.3}
\]
To justify this identification rather than assume it at
each stage, use a further universe for presheaf categories.
Restriction $q^*$ is fully faithful because $q$ is a
localization. Its left adjoint, left Kan extension along
$q^{\op}$, sends the representable
$\Hom_{\mathcal H}(-,K_i)$ to the representable
$\Hom_{\Hot}(-,BK_i)$, as the Yoneda lemma verifies on
every test presheaf. This left adjoint preserves colimits.
Its unit is invertible on a presheaf already factoring
through $q$. Applying these facts to (B.12.1) proves
(B.12.3) with its natural comparison map.

Choose a small infinite cardinal $\mu$ bounding the
cardinality of the arrows of every $K_i$. Such a bound
exists because the index category and all its categories
are small in the same universe. Let
$Y=\coprod_{c\in\pi_0Y}Y_c$ be any small CW space,
choose a base point in each component, and apply B.11 to
obtain acyclic maps $t_c:BG_c\to Y_c$. Put
$T=\coprod_cBG_c$ and $t=\coprod_ct_c$. Write
$\tau_1Y=\coprod_cB\pi_1Y_c$ and let
$v:Y\to\tau_1Y$ be its canonical first truncation in
$\Hot$. The composite $vt:T\to\tau_1Y$ is represented
by the quotient homomorphisms
$G_c\twoheadrightarrow\pi_1Y_c$.

For every $i$, restriction along that quotient gives an
\emph{isomorphism of ordinary functor categories}
\[
 \Fun\left(\coprod_cB\pi_1Y_c,K_i\right)
 \xrightarrow{\ \cong\ }
 \Fun\left(\coprod_cBG_c,K_i\right).
 \tag{B.12.4}
\]
A functor from $BG_c$ chooses an object $k$ and a
homomorphism $G_c\to\Aut_{K_i}(k)$. The target group
has cardinality at most $\mu$, so this homomorphism kills
$\ker\pi_1(t_c)$ by B.11 and factors uniquely through
$\pi_1Y_c$. A natural transformation is a family of
arrows of $K_i$ commuting with the group actions. Since
the quotient is onto, commuting with all $G_c$ actions
is equivalent to commuting with all $\pi_1Y_c$ actions.
Thus objects and morphisms factor uniquely, proving
(B.12.4) even with noninvertible endomorphisms in $K_i$
and with infinitely many components. We take components
only after proving this actual isomorphism; no invalid
commutation of $\pi_0$ with an infinite product is used.

Using (B.12.3), the resulting bijection factors as
\[
 H'(\tau_1Y)\xrightarrow{H'(v)}H'(Y)
              \xrightarrow{H'(t)}H'(T).
 \tag{B.12.5}
\]
The second map is injective. Indeed B.9 says that
$[Y,BK_i]\to[T,BK_i]$ is injective for every $i$.
Filtered colimits of compatible injections of sets are
injective: an equality in the target colimit holds at
some common stage, where injectivity applies. The
composite in (B.12.5) is bijective by (B.12.4). If
$b\circ a$ is bijective and $b$ is injective, then
$a$ is bijective: injectivity is immediate, and for
$y$ choose $x$ with $ba(x)=b(y)$ and cancel $b$.
Thus $H'(v)$ is bijective. For $Y=BX$, the checked
CATHOT C.1--C.4 identify this comparison with
$H_K(\Pi_1X)\to H_K(X)$ induced by the specified unit.
It follows that $H_K$ inverts every $\W_1$ map, since
$\Pi_1$ sends those maps to equivalences of groupoids.

We finish by checking that the conclusion is the
canonical ind-isomorphism (B.12.2). The supplied
reflector $\Pi_1:\mathcal H\to\mathcal G$, where
$\mathcal G$ is the full reflected subcategory of
groupoids, induces an adjunction on small ind-categories.
This follows directly from
\[
 \Hom_{\operatorname{Ind}(\mathcal H)}((A_i),(B_j))
   =\lim_i\colim_j\Hom_{\mathcal H}(A_i,B_j)
 \tag{B.12.6}
\]
and the reflector adjunction on its individual Hom sets.
Write $L(K)=\operatorname{Ind}(\Pi_1)(K)$ and
$\eta:K\to L(K)$ for its unit. The property just proved
is orthogonality of $K$ to every $X\to\Pi_1X$.
Taking the limit over $i$ of these bijections produces
a morphism $r:L(K)\to K$ with $r\eta=\id_K$.
The object $L(K)$ has the same orthogonality property:
each target $\Pi_1K_j$ does by reflection, and filtered
colimits preserve the resulting bijections of sets.
Consequently precomposition with $\eta$ is injective
on maps $L(K)\to L(K)$. Since
$(\eta r)\eta=\eta=\id_{L(K)}\eta$, it follows that
$\eta r=\id_{L(K)}$. This proves (B.12.2).

If $Y$ is empty, (B.12.4) is the isomorphism between
two terminal functor categories, so the argument includes
that case. No uniform bound on the sizes of arbitrary
test categories $X$ was used; only the one bound on the
small family $(K_i)$ was used. An index family unbounded
in the original universe lacks this bound and is not
covered by the statement. Thus the large-index examples
retained in the CATHOT discussion do not contradict it.
\end{proof}

The construction in B.11 follows the simplicial-cone
strategy of Kan--Thurston proofs, but the cardinal and
subcomplex arguments are supplied above. We searched
for Maunder's 1981 proof, for cardinal refinements of
Kan--Thurston, and for acyclic simple envelopes. We
verified the HNN and mitosis sources used in B.10 and
the acyclic-map criterion used in B.9--B.11; we did not
obtain Maunder's full text. Consequently no assertion
about what his proof establishes is used as a premise,
and no priority claim is attached to B.11 or B.12.
\begin{claim}{B.13}{Open}{Accessibility with only the abstract derivator axioms}
For an arbitrary derivator
$\D:\Cat_\U^{\mathrm{coop}}\to\mathrm{CAT}_\V$
satisfying Der~1--4, must $\Wlc$ be generated by a
$\U$-small set of functors as a fundamental localizer?
Must its full subcategory of $\Cat_\U^{[1]}$ satisfy
the literal Loc~(9)? We likewise leave the corresponding
accessibility assertion for $\Wc$ and the literal
accessibility assertion for $\Wone$ unresolved in this
full generality. B.4 answers the first two questions
positively for every $\U$-presentable enhanced
coefficient category, and answers all three for each
concrete coefficient theory in Part C.
\end{claim}
The obstruction to extending B.4 is the missing uniform
small family of tests or a uniform bound on witnesses.
Der~1--4 and B.3 give Kan extensions and relative descent,
but do not in our proof supply such a generating family.
For the constant and locally constant classes, the first
untreated scope is an abstract derivator with no supplied
presentable enhancement or small homotopical generating
family. For $\Wone$, accessibility is also unproved here
for a general presentable coefficient category outside the
concrete families of Part C: the accessible copower argument
for equivalences has not been extended here to its
fully faithful arrows. No explicit counterexample is supplied
to either assertion. Under
Vop\v enka's principle, the set-generation questions for
$\Wlc$ and $\Wc$ are positive by
\citet[Scholie~4.2.17]{Cisinski2006}. We neither assume
that principle nor infer an unconditional counterexample
from the lack of a proof. The exact unresolved step is
to produce a single $\lambda\in\U$ and a $\U$-small
family of arrows generating the stated classes; merely
knowing their filtered-colimit closure does not do so.

\begin{claim}{B.14}{Open}{The remaining general fibration questions}
For an arbitrary Der~1--4 derivator without an enhancement, we leave
unresolved whether every Grothendieck fibration whose transports
belong to $\Wone$ universally preserves $\Wone$ under strict
pullback, and whether every small functor factors as a member of
$\Wone$ followed by such a universally $\Wone$-preserving map.
These are respectively literal Loc~(7a) and the factorization
assertion of Loc~(7F); Loc~(7c) is proved in B.7.
For the accessible localizer $W^{\mathrm{lc}}_{L_E\Sp}$ of a general
nonzero spectrum $E$, we also leave its properness, equivalently
its Loc~(7F), unevaluated beyond the cases proved in B.7.
\end{claim}
The first two gaps are exactly the steps isolated after B.7:
the direct-image argument gives cohomological invariance, while
the identification with full faithfulness and the homotopy-fibre
criterion have only been proved here with the enhancement of C.2.
The first untreated coefficient theories have no specified such
enhancement; no counterexample derivator is supplied. For the
spectrum question the exact necessary and sufficient condition is
(B.7.7), namely generation by the aspherical projections already
in the class. The loop-homology test of C.11 decides membership
but does not prove that generation statement. This evaluation is
not supplied here already for $E=K(1)$ at a fixed prime. We do
not claim that these residual questions are open in the literature.


